\section{Statistical hypotheses and model compatibility}

A confidence interval reports a range of parameter values compatible with an
observed statistic. A hypothesis test begins with one proposed value and asks
whether the observed statistic is compatible with it.

The null hypothesis has the form
\[
H_0:\theta=\theta_0.
\]
It determines the sampling distribution used for comparison.

\subsection*{The alternative hypothesis}

A test must also specify which departures from \(\theta_0\) are relevant. Common
alternatives are
\[
H_1:\theta\neq\theta_0,
\]
\[
H_1:\theta>\theta_0,
\qquad\text{or}\qquad
H_1:\theta<\theta_0.
\]
The alternative determines which values of a test statistic count as more
extreme than the observed value. A two-sided alternative treats large deviations
in either direction as evidence. A one-sided alternative uses only the direction
specified in advance.

\begin{remarkbox}
The direction of the alternative should be chosen from the scientific question,
not after the data have been inspected.
\end{remarkbox}

\subsection*{The null distribution}

Let \(T\) be a statistic whose distribution depends on \(\theta\). Under
\(H_0\), the model gives a particular sampling distribution for \(T\). After
observation, we calculate \(t_{\mathrm{obs}}\) and locate it in that distribution.
The reasoning is
\[
H_0
\longrightarrow
\text{null sampling distribution}
\longrightarrow
 t_{\mathrm{obs}}
\longrightarrow
\text{assessment of compatibility}.
\]
An unusual observation is evidence against the null model, not a logical
contradiction.

\subsection*{A standardized statistic}

When a normal approximation is appropriate,
\[
Z=\frac{T-\theta_0}{\operatorname{SE}_{H_0}(T)}
\]
is approximately standard normal under \(H_0\). The observed value is
\[
z_{\mathrm{obs}}
=
\frac{t_{\mathrm{obs}}-\theta_0}{\operatorname{SE}_{H_0}(T)}.
\]
The sign gives the direction of the difference, and the magnitude gives the
distance from the null centre in standard-error units.

\subsection*{The p-value}

\begin{definitionbox}
The \textbf{p-value} is the probability, calculated under \(H_0\), of obtaining a
test statistic at least as extreme as the observed one in the direction or
directions specified by \(H_1\).
\end{definitionbox}

For \(Z\sim N(0,1)\), the three common forms are
\[
P(|Z|\geq|z_{\mathrm{obs}}|)
\quad\text{for }H_1:\theta\neq\theta_0,
\]
\[
P(Z\geq z_{\mathrm{obs}})
\quad\text{for }H_1:\theta>\theta_0,
\]
and
\[
P(Z\leq z_{\mathrm{obs}})
\quad\text{for }H_1:\theta<\theta_0.
\]

A small p-value means that the observed statistic, or a value more extreme in the
prespecified direction, would be unusual under \(H_0\). A large p-value means
that the observation is not unusual under \(H_0\); it does not prove that the
null hypothesis is true.

In particular,
\[
p\text{-value}\neq P(H_0\text{ is true}\mid\text{data}).
\]

\subsection*{Example: a population proportion}

For independent Bernoulli observations, consider
\[
H_0:p=p_0.
\]
Under the null hypothesis,
\[
\operatorname{SE}_{H_0}(\widehat p)
=
\sqrt{\frac{p_0(1-p_0)}{n}},
\]
so
\[
Z=
\frac{\widehat p-p_0}{\sqrt{p_0(1-p_0)/n}}.
\]
For a two-sided question, the alternative is
\[
H_1:p\neq p_0
\]
and the p-value uses both tails. For an upper-sided or lower-sided question, only
the corresponding tail is used.

\subsection*{Decision rules and possible errors}

A procedure may specify a significance level \(\alpha\). A conventional rule is
\[
p\text{-value}\leq\alpha
\quad\Longrightarrow\quad
\text{reject }H_0.
\]
Otherwise the correct statement is that the data do not provide sufficient
evidence to reject \(H_0\).

A type I error is rejection of a true null hypothesis. A type II error is failure
to reject a false null hypothesis. The significance level controls the type I
error probability under the null model. The type II error probability depends on
the particular alternative and the sample size.

\subsection*{Statistical and practical importance}

A large sample may make a small difference detectable. Statistical significance
therefore does not imply practical importance. A complete analysis should report
the estimated effect and its uncertainty, not only a threshold decision.

\subsection*{Relation to confidence intervals}

For corresponding two-sided large-sample procedures, a value \(\theta_0\) is
excluded by an approximate \(95\%\) confidence interval exactly when the
corresponding level-\(0.05\) test regards it as incompatible with the data.
One-sided tests correspond to one-sided confidence bounds rather than ordinary
two-sided intervals.

\subsection*{Dependence on the model}

A small p-value can reflect an incorrect parameter value or failure of another
assumption used to construct the null distribution. Dependence, biased sampling,
measurement error, or an inappropriate probability model can all affect the
conclusion.

\begin{summarybox}
A hypothesis test compares an observed statistic with the distribution predicted
by \(H_0\). The alternative hypothesis determines the relevant tail or tails, so
it is part of the definition of extremeness and of the p-value.
\end{summarybox}
